\documentclass[12pt]{article}
\usepackage[utf8]{inputenc}
\usepackage[T1]{fontenc}
\usepackage{hyperref}
\title{Zay Language Resources for K-12 Curriculum Development}
\author{Compiled by Claude Sonnet 5.5 (claude-sonnet-5-5) for Andr\'as Kornai}
\date{}
\begin{document}
\maketitle

\section{Introduction}
Zay (Lak'i, Laqi; ISO 639-3 zwa; Glottocode zayy1238) is an Ethiopian Semitic language spoken around Lake Zway. Sources give speaker numbers from about 4,900 to 14,000. Use of Zay is declining as speakers shift to Oromo. It is not known to be taught in schools, and one source describes it as unwritten. Little material was found.

\section{Primary Educational Resources}
No primers, textbooks or school materials were found.

\section{Linguistic and Reference Materials}
\begin{itemize}
\item Wikipedia entry on the Zay language, whose reference list cites Ronny Meyer (2002) on Oromo influence on Zay vocabulary, Getu Shikur (1999) on Zay morphology, and Wolf Leslau (1999), \textit{Zway Ethiopic Documents}. \url{https://en.wikipedia.org/wiki/Zay_language}
\item Wikipedia entry on the Zay people. \url{https://en.wikipedia.org/wiki/Zay_people}
\item Catalogue of Endangered Languages entry for Zay. \url{https://endangeredlanguages.com/node/88388}
\end{itemize}

\section{Religious and Community Resources}
\begin{itemize}
\item Scripture Earth page for Zay. It points to audio recordings on Bible.is and from Global Recordings Network, and no written text was found. \url{https://www.scriptureearth.org/00eng.php?iso=zwa}
\item Joshua Project profile of the Zay people in Ethiopia. \url{https://joshuaproject.net/index.php/people_groups/16094}
\end{itemize}

\section{Digital and Audio Resources}
Audio religious recordings are listed on the Scripture Earth page above.

\section{Curriculum Development Context}
No dictionary, written scripture or school programme was confirmed. Any curriculum would need new documentation work with the Zay community. The existing linguistic studies and recordings could serve as a starting point.

\end{document}